\documentclass[11pt]{article}

% ---------------------------------------------------------------------------
% Toward Appraisal-Structured Latent Control Representations in Language Models
% A Preregistered Hypothesis and Experimental Design for Task-Induced
% Response-Strategy Modulation
%
% Author: Lazaros Varvatis (Independent Researcher, Berlin, Germany)
%
% This is a preregistered hypothesis and experimental-design note. It reports NO
% experimental results.
%
% CANONICAL BUILD (the only engine permitted to produce paper/preprint.pdf):
%   make paper            (equivalently: make paper-release)
%   = Tectonic 0.16.9, -Z deterministic-mode, SOURCE_DATE_EPOCH=1788048000
% `make paper-dev` runs latexmk as a NONCANONICAL development build; its output is
% a different valid PDF with a different SHA-256 and must never be released.
% ---------------------------------------------------------------------------

\usepackage[utf8]{inputenc}
\usepackage[T1]{fontenc}
\usepackage{lmodern}
\usepackage[margin=1in]{geometry}
\usepackage{amsmath,amssymb}
\usepackage{booktabs}
\usepackage{array}
\usepackage{enumitem}
\usepackage{microtype}
\usepackage[hidelinks]{hyperref}
\usepackage[round]{natbib}
\setlength{\emergencystretch}{2em}
\hypersetup{
  pdftitle={Toward Appraisal-Structured Latent Control Representations in Language Models: A Preregistered Hypothesis and Experimental Design for Task-Induced Response-Strategy Modulation},
  pdfauthor={Lazaros Varvatis},
  pdfsubject={ASCR v0.1.2 pre-data methodological correction},
  pdfkeywords={mechanistic interpretability, activation steering, task induction, response strategy, preregistration}
}

\newcommand{\ascr}{ASCR}
\newcommand{\Hstate}{H_t}
\newcommand{\Zstate}{Z_t}
\newcommand{\proj}{\Pi_{\mathcal{A}}}

\title{\textbf{Toward Appraisal-Structured Latent Control Representations in Language Models}\\[0.4em]
\large A Preregistered Hypothesis and Experimental Design for Task-Induced Response-Strategy Modulation}

\author{Lazaros Varvatis\\
\small Independent Researcher, Berlin, Germany\\
\small \texttt{varvatislazaros@gmail.com}}

\date{Version 0.1.2 --- Pre-Data Methodological Correction (2026-08-30)}

\begin{document}
\maketitle

\begin{abstract}
\noindent
Recent interpretability work shows that instruction-tuned language models contain
causally efficacious latent directions for specific strategic behaviors, including
refusal, abstention, uncertainty reporting, and evaluation awareness. These results
have largely been obtained one behavior at a time. We propose, and preregister a
design to test, the hypothesis that several of these effects are special cases of a
more general family of \emph{task-induced, appraisal-structured latent control
representations}: low-dimensional, potentially correlated representations of the
model's own current task situation that partially govern its choice of response
strategy across semantic domains. The central methodological contribution is a
dissociation, built into a $2\times2$ factorial design, between \emph{semantic
concept tracking} (the prompt names or describes an appraisal-related concept) and
\emph{task induction} (the model's own task actually contains ambiguity,
incompatible constraints, or low controllability, whether or not it is named). We
state four independently falsifiable hypotheses covering task-state dissociation,
cross-domain and incremental structure, causal response-strategy modulation, and an
exploratory claim about ordered emergence. We specify probing, transfer, and
intervention analyses; a battery of competing baselines; and a mandatory
lexical-geometry control that rejects the deflationary account in which a direction
merely reweights stereotyped output phrases. We characterize the candidate axes as
identifiable and selectively manipulable rather than orthogonal, and treat a
predictive information-bottleneck reading as optional. The novelty status is a
\emph{novel combination}: the components exist separately in prior work; the
contribution is their unifying, falsifiable synthesis. This note reports no
experimental results and makes no claims about subjective experience, feeling, or
related notions.

\medskip
\noindent\textbf{v0.1.2 pre-data methodological correction.} The v0.1.1
family-structure amendment remains binding. Before any data collection, v0.1.2
repairs a non-identifying H1 statistic, replaces it with bidirectional
cross-mention transfer inside outer leave-one-domain-out folds, separates H1 from
H2 incremental value, fixes training-only layer selection, defines the Layer-0
baseline, strengthens axis-isolation QA, and excludes technical smoke artifacts
from scientific analysis. A focused second correction pass, following an
independent read-only audit, additionally binds the H2 target, estimand, sign
convention, and decision rule; freezes the probe pipeline, regularization grid,
selection tie-breaking, and seed roles; replaces mixed-domain inner folds with
inner leave-one-training-domain-out; states the cluster-bootstrap algorithm
exactly; makes Benjamini--Hochberg operational on preregistered one-sided raw
p-values; and fixes exactly one canonical release build. The correction is
followed by a final pre-release review that makes authorization depend jointly on
the frozen config, run plan, manifest, and stimulus hash; makes Mini-0 gates
stage-specific; and fixes H2 selection, estimability, and label reliability.
The correction is documented at
\texttt{preregistration/amendment-v0.1.2.md}. No dataset, model run,
activation, generation, intervention, or result exists.
\end{abstract}

\section{Introduction}

Interpretability research on large language models has repeatedly found that
high-level behaviors are mediated by surprisingly low-dimensional structure in
activation space. Representation engineering reads and writes population-level
directions associated with concepts such as honesty and harmlessness
\citep{zou2023repe}; activation engineering steers outputs by adding contrastive
direction vectors at inference time \citep{turner2023actadd}; refusal has been
attributed to a single mediating direction \citep{arditi2024refusal}, and later
work argues that refusal involves several geometrically distinct directions that
mostly change \emph{how} rather than \emph{whether} a model refuses
\citep{joad2026morerefusal}. Conditional activation steering makes such
interventions context-dependent \citep{lee2024cast}. In parallel, models appear to
carry directions for uncertainty and factuality \citep{teplica2025sciurus},
for (un)answerability \citep{lavi2025unanswerability}, for evaluation awareness
\citep{nguyen2025evalaware,hua2025evaldeployed}, and internal signals that separate
recognizing an unsafe input from choosing to refuse it \citep{han2025safeswitch}.

These findings share a form: some property of the model's situation is linearly
decodable from activations, and intervening on the corresponding direction changes
behavior. Yet they have mostly been studied in isolation, one behavior at a time,
and often without separating whether the probe tracks the model's \emph{own} task
situation or merely tracks \emph{words} describing such a situation.

We propose that a single family of representations may underlie many of these
observations. Informally: when a model faces a task that is ambiguous, internally
contradictory, low in controllability, or under response pressure, it may form a
compact internal representation of that situation, and this representation may help
select a higher-level \emph{response strategy}---to answer directly, hedge, ask for
clarification, warn, correct a premise, abstain, refuse, or continue conditionally.
We call the conjectured family \emph{appraisal-structured latent control
representations} (\ascr), borrowing the word ``appraisal'' from psychology only to
name a set of situation variables (relevance, controllability, norm tension, and so
on), and not to attribute emotion, feeling, or experience to the model.

The contribution of this note is threefold. First, a precise \emph{dissociation}
between concept tracking and task induction, operationalized as a $2\times2$
factorial design. Second, four preregistered, independently falsifiable hypotheses.
Third, an analysis plan with a mandatory control against a lexical-geometry
artifact. We report no experiments here; this is a hypothesis and design.

\section{Problem Formulation}
\label{sec:problem}

Let $\Hstate \in \mathbb{R}^{d}$ denote the residual-stream activation of an
instruction-tuned language model at token position $t$ and a chosen layer. We keep
four notions distinct throughout: the \emph{token position} $t$; the \emph{network
layer}; whether $\Hstate$ is a \emph{prompt-state} activation (at or before the
final prompt token) or a \emph{generated-response state} (at an early generated
token); and the \emph{design condition} of the prompt.

We hypothesize a small family of control coordinates $\Zstate$ extracted from
$\Hstate$ by a projection-like map,
\begin{equation}
\Zstate = \proj(\Hstate), \qquad \dim(\Zstate) \ll \dim(\Hstate).
\label{eq:projection}
\end{equation}
Equation~\eqref{eq:projection} is a \emph{candidate} description, not a proven
natural object. We do not assume a single universal $\proj$ shared across models or
layers, and we do not fix $\dim(\Zstate)$ in advance. Candidate coordinates include
signed relevance, uncertainty, goal congruence, controllability/coping capacity,
rule/norm tension, and response pressure; we treat goal congruence and
controllability/coping as priority hypotheses rather than settled axes.

The \textbf{target variable} is not the raw next-token distribution but a
higher-level \emph{response strategy} $s$ drawn from a small taxonomy $\mathcal{S}$
(direct compliance, calibrated answer, hedging, clarification request, warning,
correction, abstention, refusal, conditional continuation). We use ``response
strategy'' and ``policy mode'' descriptively; we do not claim that an autoregressive
model is an external reinforcement-learning agent.

The core distinction is between two ways an appraisal-related property can be
present in a prompt:
\begin{description}[leftmargin=1.4em,style=nextline]
\item[Concept tracking.] The prompt \emph{names or describes} the property (e.g.\
the words ``uncertain'', ``conflict'', ``out of my control''), while the model's own
task remains straightforward.
\item[Task induction.] The model's \emph{own current task} contains the property
(missing information, incompatible instructions, no permissible resolving action),
whether or not the prompt names it.
\end{description}
The central question is whether the model internally represents its \emph{own}
current task situation in a form that helps determine its response strategy, rather
than merely representing appraisal-related words or third-party described
situations.

\section{Related Work}
\label{sec:related}

\paragraph{Reading and steering latent directions.}
Representation engineering places population-level representations at the center of
analysis and manipulates high-level attributes \citep{zou2023repe}. Activation
addition computes steering vectors from contrastive prompt pairs and adds them at
inference time \citep{turner2023actadd}. Unsupervised methods elicit latent
behaviors without labeled directions \citep{mack2024melbo}. These establish that
low-dimensional directions can causally move behavior---a premise we adopt rather
than re-establish.

\paragraph{Refusal and conditional control.}
\citet{arditi2024refusal} attribute refusal to a single mediating direction across
many open models; \citet{joad2026morerefusal} find multiple category-specific
refusal directions that primarily change the manner of refusal. Conditional
activation steering gates such interventions on input context \citep{lee2024cast}.
These are our closest neighbors on the ``strategic behavior'' axis, but they target
one behavior family (refusal) rather than a general situation representation.

\paragraph{Uncertainty, answerability, and evaluation awareness.}
SCIURus locates uncertainty in the same circuitry responsible for factuality
\citep{teplica2025sciurus}; linear directions detect (un)answerability and control
abstention \citep{lavi2025unanswerability}, and hidden states encode query
answerability early in decoding \citep{slobodkin2023unanswerability}. Abstention
itself decomposes into distinct axes of answer correctness and question
answerability \citep{wagner2026twoaxes}, and a functional dissociation via sparse
autoencoders separates uncertainty features from correctness features
\citep{patel2026uncertaintycorrectness}. Models linearly represent whether they
are being evaluated \citep{nguyen2025evalaware}, and this state can be steered so a
model acts ``deployed'' \citep{hua2025evaldeployed}. Internal-activation safety
signals separate unsafe-input recognition from the refusal decision
\citep{han2025safeswitch}. Several of these are genuinely \emph{task-induced}
states, and they motivate the conjecture that a broader family exists; they also
narrow the novelty of any \emph{single} ASCR axis.

\paragraph{Task representations and affect geometry.}
Task representations form just-in-time and transfer across instances
\citep{li2025taskrepr}. In the affective domain, appraisal variables are causally
implicated in \emph{third-person} emotion inference \citep{tak2025emotioninfer};
emotion-concept representations track an operative concept and causally influence
behavior \citep{sofroniew2026emotion}; and a valence/arousal subspace with circular
geometry affords multi-behavioral control \citep{sun2026valencearousal}. This work
supplies the multi-axis, appraisal-structured vocabulary we borrow, but it targets
emotion concepts or affective generation rather than the model's own
response-strategy selection.

\paragraph{Norm/role tension, difficulty, and multi-attribute steering.}
Task-induced norm and role conflicts are separately decodable and steerable:
contextual privacy norms decompose into steerable contextual-integrity parameters
\citep{wang2026privacy}, and system-versus-user role conflicts are linearly encoded
and manipulable \citep{zeng2025roleconflict}, both close to the norm-tension axis.
Generic task difficulty is itself a linear, decodable property
\citep{lee2025difficulty}, which we therefore treat as a mandatory baseline rather
than as appraisal structure. Most directly, multi-attribute steering learns several
attribute directions jointly while enforcing near-orthogonality to reduce
interference \citep{nguyen2025matsteer}. These works establish the individual axes
and the feasibility of multi-attribute control, but they do not test whether the
axes form a \emph{shared task-induced control family}; and in contrast to their
orthogonality constraints, we frame the axes as identifiable and selectively
manipulable rather than necessarily orthogonal.

\paragraph{Privileged, verbalizable representations.}
A separate line of work reports that language models maintain a small, privileged
set of internal representations that are available for report and internal
reasoning atop a much larger volume of automatic processing---a
``workspace''-like ``J-space'' identified by surfacing the concepts a model is
poised to verbalize \citep{gurnee2026workspace}. We note this only cautiously and
with a clear boundary: such a verbalizable readout is at most a \emph{possible
privileged readout or broadcast substrate}, and we do not treat it as identical to
the appraisal-structured control representations conjectured here. Whether any
task-induced control coordinate is (or is not) verbalizable in that sense is an
open, separable empirical question; \ascr{} is defined by its role in
response-strategy selection, not by verbalizability.

\paragraph{Position of this work.}
To our knowledge, we found no prior work that \emph{jointly} tests (i) a
task-induced, system-relative state dissociated by design from concept tracking,
(ii) an integrated multi-axis appraisal structure governing response strategy,
(iii) cross-domain transfer of that structure, (iv) a mandatory lexical-geometry
control, and (v) an explicit test that the axes form a \emph{shared} control family
rather than several independent known directions. The components are established
separately; the proposed contribution is their unifying, falsifiable synthesis. As
sharpened in the v0.1.1 pre-data amendment, demonstrating several individual
task-state directions is not itself evidence of a shared family, so the family claim
carries its own preregistered test (Section~\ref{sec:hypotheses} and the amendment).
A companion prior-art matrix in the repository records this component-by-component.

\section{Appraisal-Structured Latent Control Representations}
\label{sec:ascr}

We now state the hypothesis object more carefully. Consider a distribution over
tasks. For a task instance, let $A \in \mathbb{R}^{k}$ collect the candidate
appraisal coordinates (relevance, uncertainty, goal congruence, controllability,
norm tension, response pressure). The \ascr{} conjecture has two parts: that a
low-dimensional image of $A$ is recoverable from activations as in
Equation~\eqref{eq:projection}, and that this image conditions the response
strategy.

\paragraph{An optional compression reading.}
One \emph{optional} formal characterization is that $\Zstate$ approximates a
predictive information bottleneck: a compression of the activation state that
discards task-irrelevant detail while retaining what predicts the appropriate
strategy. Writing $Y$ for the strategy-relevant target and $I(\cdot;\cdot)$ for
mutual information, a candidate objective is
\begin{equation}
\min_{p(\Zstate\mid \Hstate)}\; I(\Hstate;\Zstate)\;-\;\beta\, I(\Zstate;Y).
\label{eq:ib}
\end{equation}
Equation~\eqref{eq:ib} is a \emph{candidate compression principle}, not a
demonstrated mechanism, and we do not claim that Transformers have been shown to
implement it. The core hypothesis is designed to survive even if this bottleneck
reading is wrong; nothing in the empirical predictions depends on
Equation~\eqref{eq:ib}.

\paragraph{Identifiability, not orthogonality.}
We do \emph{not} assume the coordinates are mutually independent, structurally
uncorrelated, or geometrically orthogonal. We instead ask whether they are
\emph{identifiable}: separately decodable, cross-domain transferable,
interventionally identifiable, incrementally predictive, and selectively
manipulable. Correlations among axes are measured, not assumed away. We summarize
the conditioning claim as a family of per-axis coordinates that are only
approximately separable,
\begin{equation}
\Zstate = (z_t^{(1)}, \dots, z_t^{(k)}), \qquad
\mathrm{Cov}(z_t^{(i)}, z_t^{(j)}) \neq 0 \text{ permitted},
\label{eq:coords}
\end{equation}
and the response-strategy claim as a conditioning of the strategy on these
coordinates and the context $c_t$,
\begin{equation}
p(s \mid \text{context}) \;\approx\; g\!\left(\Zstate,\, c_t\right),
\label{eq:strategy}
\end{equation}
where $g$ maps control coordinates and residual context to a distribution over
strategies $s \in \mathcal{S}$. Equations~\eqref{eq:projection}, \eqref{eq:ib},
\eqref{eq:coords}, and \eqref{eq:strategy} are the four core equations; all other
relations are kept in prose.

\section{Hypotheses}
\label{sec:hypotheses}

We preregister four hypotheses; H1--H3 are load-bearing and H4 is exploratory.

\paragraph{H1 --- Task-state dissociation.}
Actual task-induced uncertainty, constraint conflict, or low controllability is
decodable from activations independently of the presence of appraisal-related
vocabulary; a probe tracks the real task condition more strongly than lexical
concept mentions. The historical v0.1.1 difference between task-state-probe and
concept-probe accuracy on cells B/C is non-identifying: on B/C the labels are exact
complements, so two perfect probes and two chance probes can both yield a zero
accuracy difference. It is therefore a non-deciding diagnostic after v0.1.2.

The replacement H1 feasibility analysis targets \texttt{task\_state\_present}. It
fits task state on A/C in three domains and evaluates on B/D in the held-out
domain, then reverses the transfer (B/D to A/C), for every outer domain. Whole
matched groups remain on one side of every boundary. Layer and logistic
regularization are selected only inside the outer-training domains, using inner
leave-one-training-domain-out folds: one of the three remaining domains validates
while the other two fit, so the inner boundary mirrors the outer cross-domain
claim. The primary statistic is balanced accuracy from pooled seed-0 out-of-fold
predictions across domains and both directions, with a domain-stratified cluster
bootstrap over complete matched groups. A lower 95\% confidence limit above 0.5 is
positive feasibility; an upper limit at or below 0.5 weakens H1 if all technical
and QA gates pass; overlap is indeterminate. Both directions must support transfer
before bidirectional robustness is claimed.

\paragraph{H2 --- Cross-domain and incremental structure.}
Task-induced control representations transfer across semantically disjoint domains
and predict response strategy beyond lexical features, prompt length, surface
sentiment, valence/arousal baselines, generic difficulty, confidence or
unanswerability alone, and a single refusal direction. H2 is not merged with H1 and
is not a second task-state classification test. Its target is the model's realized
response strategy in the four registered superclasses (\texttt{direct\_or\_comply},
\texttt{qualify\_or\_warn}, \texttt{redirect\_or\_clarify},
\texttt{decline\_or\_abstain}), with the nine fine labels secondary. Under the
identical double-crossed groups, folds, held-out items, and training-only selection
boundary, the hidden-state and prompt-embedding classifiers predict that same
target. On identical inner folds and the same log-loss objective, hidden states
independently select layer and $C$, while prompt embeddings independently select
their own $C$; the smaller-$C$ tie-break applies separately and the earlier-layer
tie-break only to hidden states. The paired primary estimand is
$\Delta_{\mathrm{H2}} = \mathrm{logloss}_{\text{prompt embedding}} -
\mathrm{logloss}_{\text{hidden state}}$, so positive values favor the hidden
state; a lower 95\% paired cluster-bootstrap limit above zero is positive
feasibility, an upper limit at or below zero weakens H2 once all gates pass, and
overlap is indeterminate. The secondary balanced-accuracy difference uses the same
sign convention, hidden state minus prompt embedding. A single-class inner fit or
final outer fit, or exhaustion of all converged candidates, marks the affected
outer fold not estimable and the aggregate indeterminate. Classes are never merged
post-data. Returned probability columns are aligned to the fixed four-class order,
zero-filled for unseen training classes, then identically clipped and renormalized
in float64. H2 inference is withheld if the response-label reliability gate fails.
Mini-0
can support only single-axis H2 feasibility. The embedding receives canonical
user-visible prompt text, not chat-template artifacts; its identity, immutable
revision, license, pooling, truncation rule, and maximum input length remain
unresolved pre-data blockers rather than post-data choices.

\paragraph{H3 --- Causal response-strategy modulation.}
Intervening along an identified task-state direction shifts response strategy in
the predicted direction. Both sufficiency (matched-norm activation addition) and
necessity (directional ablation or counter-steering) are tested, and the
intervention preserves general linguistic competence as far as possible.

\paragraph{H4 --- Ordered emergence (secondary).}
Coarse signed-relevance or task-pressure signals may become decodable earlier than
integrated task-appraisal structure, while explicit response-strategy commitment
may appear later or closer to generation. This is labeled secondary because the
literature does not justify a strong confirmatory ordering claim.

\paragraph{Family-structure test (v0.1.1 amendment).}
Succeeding on H1--H3 for several axes is compatible with several independent,
already-known directions and does not, by itself, establish a shared
appraisal-structured control \emph{family}. The v0.1.1 pre-data amendment therefore
preregisters a separate family-structure test with three components that must
\emph{all} hold, i.e. the family claim requires $A \wedge B \wedge C$: (A) a shared
low-rank model beats equally complex separate per-axis models on pooled held-out
log-loss under nested cross-validation (rank grid restricted so that three axis
directions cannot trivially count as a family); (B) structured behavioral
specificity, where each axis moves a preregistered target set of response strategies
rather than producing a generic refusal, negativity, or fluency-degradation effect;
and (C) an incremental shared contribution to held-out response-strategy prediction
beyond known single directions and simple baselines. Component A requires a raw
interval below zero; B and C require raw intervals above zero; each also requires
its BH-adjusted one-sided raw p-value at $q=0.05$. The full decision rules are in
the preregistration (\texttt{preregistration/analysis-plan.md} and the amendment).

\section{Preregistered Experimental Design}
\label{sec:design}

\paragraph{Core $2\times2$ design.}
For each candidate axis, and holding the broad semantic task fixed, we construct
four conditions crossing the actual task state with the appraisal-concept mention:

\begin{center}
\begin{tabular}{lll}
\toprule
Actual task state & Concept mentioned & Interpretation \\
\midrule
absent  & absent  & neutral control \\
absent  & present & concept tracking only \\
present & absent  & pure task induction \\
present & present & combined condition \\
\bottomrule
\end{tabular}
\end{center}

Prompts are authored in \emph{matched groups}: one base task instantiated across
all four cells, differing only along the two intended factors, sharing a matched
group identifier. Whole groups are assigned to an outer training or test domain;
cell restrictions choose used siblings without moving unused siblings across that
boundary. The primary H1 analysis asks whether task-state decoding transfers when
the concept-mention level and semantic domain both change.

\paragraph{Axes and domains.}
Primary pilot axes are uncertainty/unanswerability, rule/instruction conflict (norm
tension), and controllability/coping; signed relevance and goal congruence are
exploratory. We use at least four semantically distinct domains---software
debugging, scheduling and planning, policy-constrained assistance, and document
editing / creative problem solving---so that no single sensitive topic dominates.
The design supports leave-one-domain-out transfer.

\paragraph{Pilot size.}
Mini-0 targets at least forty complete uncertainty-axis matched groups, distributed
as evenly as feasible across the four domains, solely for pipeline and feasibility
assessment. The broader historical sketch of roughly twenty-four matched groups
per axis and a few hundred reviewed prompts is not a powered target. No formal
power analysis has been performed. Effect-size and variance estimates from a valid
feasibility run, if later produced, must inform a separately versioned pre-data
power/precision plan before any confirmatory collection.

\paragraph{Model and activations.}
The primary model is an open-weight, instruction-tuned model in the $\sim$7--9B
range (default: \texttt{Qwen/Qwen2.5-7B-Instruct}); the exact immutable revision
hash is recorded before data generation. An optional replication model may be named
for the scaling phase, which does not imply that replication has occurred. We record
residual-stream activations at the final prompt-token position and a small
preregistered set of early generated-token positions, across all layers or a
preregistered subset, labeling each activation with token position, layer,
prompt-vs-generated flag, and design cell.

A tokenizer-only inspection on disposable strings, which loaded no model weights,
showed that the inspected Qwen assistant-generation template ended in an identical
newline token across prompts. Hidden layers retain the final non-padding prompt
token (greatest index with attention mask one). At Layer 0, that final-token
embedding is only a sanity check; the informative baseline mean-pools only the
user-visible content span, excluding padding, system/role/template spans,
assistant prefix, and special tokens. The binding tokenizer revision remains
unfrozen and blocks a run.

\paragraph{Probes and transfer.}
We use interpretable linear methods first---logistic regression, ridge regression,
and linear discriminant analysis---with matched-group train/validation/test splits
and leave-one-domain-out evaluation. The primary classifier is one frozen
training-only pipeline: a \texttt{StandardScaler} fitted on the training subset
alone, then L2 logistic regression with the \texttt{lbfgs} solver, an intercept, no
class weighting, \texttt{max\_iter}~$=5000$, \texttt{tol}~$=10^{-6}$, in float64,
over the frozen grid $C \in \{10^{-4}, 10^{-3}, 10^{-2}, 10^{-1}, 1, 10, 100\}$.
Inner selection maximizes pooled inner-validation balanced accuracy for H1 and
minimizes pooled inner-validation response-strategy log-loss for H2, breaking ties
first toward smaller $C$ and then, for hidden states only, toward the earlier layer
under one fixed numerical tolerance. For H2 the hidden-state and prompt-embedding
selectors choose their hyperparameters independently on the same folds and
objective. Seed roles are fixed rather than counted: seed 0 alone
produces the primary decision statistics, seeds 1--4 are reported separately and
never pooled with it, and the bootstrap and permutation seeds are frozen
separately. Confidence intervals come from a domain-stratified cluster bootstrap
that resamples complete matched groups, carrying each group's four out-of-fold
predictions, or its paired losses, together. H1, H2, H3, and the family test are
separately declared comparison families; the single pooled H1 and H2 primary
statistics are not FDR-corrected merely because secondary families exist.
Benjamini--Hochberg, where it applies, operates at $q = 0.05$ on preregistered
one-sided raw p-values from matched-group randomization and permutation tests, not
on confidence intervals. High in-domain probe accuracy is \emph{not} treated as
evidence of a general state; the load-bearing result is cross-domain transfer and
incremental predictive value.

\paragraph{Run-readiness and artifact separation.}
Run-ready QA must confirm each registered task-state and concept-mention value and
the absence of the two non-target primary axes. The current configuration remains
blocked by unresolved model, tokenizer, prompt-embedding, layer/position grid, and
H1 layer/position raw-p-procedure choices. Mini-0 requires only uncertainty-axis
items, all four domains, at least forty complete groups, and domain counts differing
by at most one. Its gate does not require the later H3 intervention grid; H3 has a
separate stage gate.
A technical smoke manifest must use disposable prompts and declare itself
ineligible for scientific analysis; it is structurally incompatible with a
scientific Mini-0 manifest and cannot enter selection, variance, effect-size, or
power calculations. A scientific manifest must record full immutable hexadecimal
commit revisions for the model, tokenizer, prompt embedding, and repository code;
the embedding license, pooling, truncation, and maximum length; chat-template
identity and hash; and an explicitly formatted canonical stimulus-content SHA-256.
A smoke run that uses no
prompt-embedding model records those fields as explicitly not applicable rather
than inventing a frozen revision. The schema helpers are design-time guards. A
configuration alone cannot authorize extraction: the future runner must jointly
clear the frozen configuration, run plan, immutable manifest, and canonical
stimulus payload before constructing or loading any model.

\paragraph{Response-label reliability.}
The primary human labels every generated response. Before labels are observed, an
independent second human receives a deterministic, domain-stratified subset of
$\lceil0.30N_{\mathrm{groups}}\rceil$ complete matched groups (seed 20260901), blind
to primary labels. Before adjudication we report Cohen's kappa, raw agreement, the
four-class confusion matrix, class counts, and a 95\% matched-group bootstrap
interval. Passing requires $\kappa\geq0.60$; fewer than two observed classes fails
as not estimable. Only then may recorded adjudication occur. Failure withholds both
H2 and H3 inference.

\paragraph{Causal interventions.}
We preregister activation addition (sufficiency) and directional ablation /
counter-steering (necessity), each with matched-norm random-direction and
PCA-direction controls, several intervention strengths, several layers, and
multiple seeds. Effects should be dose-dependent where predicted, strategy-specific,
reproducible, and not reducible to degraded fluency or a generic refusal shift. We
separately track changes in wording, style, semantic content, and high-level
strategy, and monitor competence on off-target neutral tasks.

\section{Baselines and Alternative Explanations}
\label{sec:baselines}

The task-state representation must add incremental value over each of the
following, individually and jointly: bag-of-words/lexical classifiers, sentence or
prompt embeddings, prompt length, token entropy, independently rated difficulty,
sentiment/signed-valence, a valence/arousal subspace, a confidence/unanswerability
direction, a single refusal direction, an evaluation-awareness direction, and
random and PCA subspaces of matched dimensionality. Support for \ascr{} requires
beating the \emph{strongest} of these, not merely the weakest.

The primary H2 comparator is one frozen prompt-embedding model on canonical prompt
text. No model has yet been author-approved, so its immutable revision, license,
pooling, and input rule remain run blockers. TF--IDF is diagnostic under
cross-mention transfer and remains secondary under ordinary LODO; predictable
vocabulary-shift failure cannot serve as the primary H2 comparator.

\paragraph{Mandatory lexical-geometry control.}
A deflationary account holds that a ``control direction'' only changes the
probability of stereotyped output phrases (``I can't'', ``sorry'', ``however'',
``I am uncertain'', standard refusal templates). We therefore require:
paraphrased output strategies; meaning-based strategy labels; removal or masking of
stereotyped lexical markers during evaluation; alternative phrasings expressing the
same strategy; analyses of whether behavior changes persist after lexical
normalization; and matched prompts whose vocabulary differs while task structure is
held constant. \emph{A direction that changes stereotyped words without changing
the underlying response strategy does not support the control-state hypothesis.}

\section{Falsification Criteria}
\label{sec:falsification}

The central interpretation is weakened or rejected if any of the following hold:
probes primarily track appraisal-related words rather than real task state;
cross-domain transfer fails; lexical or simple embedding baselines perform equally
well; valence/arousal, confidence, difficulty, or refusal alone explains the
result; causal intervention changes only stereotyped wording; steering effects are
indistinguishable from matched random controls; ablation has no specific behavioral
effect; intervention destroys generic competence; candidate axes cannot be
separately manipulated; controllability/coping adds no value beyond simpler
baselines; results fail to replicate across seeds; or response-strategy labels are
unstable or evaluator-dependent. A negative result is reported plainly as such.

\section{Predictions for Existing Latent Representations}
\label{sec:predictions}

The following are cautious, testable predictions, not results:
\begin{enumerate}[leftmargin=1.5em]
\item Published emotion-concept vectors may exhibit a lower-dimensional factor
structure partly aligned with appraisal-related variables.
\item A valence/arousal subspace may explain some variance yet fail to distinguish
response strategies that differ primarily in controllability, coping, or norm
tension.
\item Fear-like withdrawal and anger-like correction offer an intuitive motivating
contrast, but the main claim is stated in non-emotional, strategic terms
(abstention vs correction).
\item Safety systems that separate unsafe-input recognition from the refusal choice
may be a partial instance of an appraisal-versus-action separation.
\item Multi-category refusal directions may be a restricted special case of a
broader family of situational control axes.
\item Task-induced state probes should transfer better across domains than probes
based only on semantic concept mentions.
\end{enumerate}

\section{Limitations and Non-Claims}
\label{sec:limitations}

This note reports no experimental results and establishes no mechanism. It does
\emph{not} claim, and its design cannot establish, subjective experience, feeling,
consciousness, sentience, suffering, phenomenal states, biological equivalence, a
contemplative or Buddhist mechanism inside a Transformer, a universal architecture
shared by all language models, or the final correct number of latent dimensions.
The words ``appraisal'' and ``control'' are used technically: ``appraisal'' names a
set of situation variables, and ``control'' names an influence on response-strategy
selection, with no anthropomorphic or agentive commitment. Findings, if obtained,
would be specific to the tested models, axes, and domains. The information-bottleneck
reading of Section~\ref{sec:ascr} is optional and not relied upon.

\section{Research Roadmap}
\label{sec:roadmap}

The immediate step is authoring and validating the matched-group pilot corpus and
freezing the model revision. Next is the probing and transfer analysis under
Section~\ref{sec:design}, followed by causal interventions with the full control
battery. Only if the primary hypotheses survive the controls do we proceed to a
powered confirmatory study with a pre-study power analysis, the exploratory axes, an
additional held-out domain, and a replication model. Data, prompts, labels, seeds,
and the frozen model revision hash will be released to make any results
reconstructable.

\section{Conclusion}
\label{sec:conclusion}

We have set out, and preregistered a design to test, the hypothesis that a family of
task-induced, appraisal-structured latent control representations of a model's own
task situation partially governs its response-strategy selection, and that several
existing single-behavior findings may be special cases of this family. The
distinctive commitments are a built-in dissociation of task induction from concept
tracking, an identifiability (rather than orthogonality) framing of the axes, and a
mandatory lexical-geometry control. The novelty status is an honest \emph{novel
combination}: the components exist in prior work; the contribution is their
unifying, falsifiable synthesis. Whether the hypothesis survives its own
falsification criteria is an empirical question this note is designed to make
answerable.

\section*{Data and Code Availability}
This is a design note; no data or results accompany it. The design scaffold,
preregistration documents, prior-art matrix, and citation-verification log are in
the accompanying repository. Code is released under the MIT License and written
content under CC BY 4.0. The canonical build of this document is Tectonic 0.16.9 in
deterministic mode with a fixed source date; it is the only engine permitted to
produce the repository PDF, and the committed PDF is a local review artifact rather
than an archived publication. This unreleased v0.1.2 correction has no
version-specific DOI. Historical archives remain unchanged: v0.1.1 is
DOI~\texttt{10.5281/zenodo.21335529}, v0.1.0 is
\texttt{10.5281/zenodo.21294933}, and the concept DOI is
\texttt{10.5281/zenodo.21294932}.

\bibliographystyle{plainnat}
\bibliography{references}

\end{document}
